%%%%%%%%%%%%%%%%%%%%%%%%%%%%%%%%%%%%%%%%%%%%%%%%%%%%%%%%%%%%%%%%%%%%%%
%% Data section
%% Paper: Mapping the land-use footprint of Brazilian soy embodied in
%%        international consumption -- a spatially explicit input-output
%%        approach
%% Target: Nature Portfolio (Nature Food) -- see docs/paper_plan.md
%%
%% Sources and links traced from code/download_data.py, code/pipeline/,
%% code/prep/, and the data/*/README.md files. All external URLs were
%% checked and resolve as of July 2026 (Brazilian government portals may
%% geo-block automated requests but resolve in a browser).
%%
%% Uses the official Springer Nature template (sn-jnl.cls) with the
%% Nature Portfolio reference style (sn-nature). Compile with:
%%   pdflatex data && bibtex data && pdflatex data && pdflatex data
%%%%%%%%%%%%%%%%%%%%%%%%%%%%%%%%%%%%%%%%%%%%%%%%%%%%%%%%%%%%%%%%%%%%%%

\documentclass[pdflatex,sn-nature]{sn-jnl}% Nature Portfolio reference style

%%%% Standard packages
\usepackage{graphicx}%
\usepackage{amsmath,amssymb,amsfonts}%
\usepackage{amsthm}%
\usepackage{xcolor}%
\usepackage{textcomp}%
\usepackage{manyfoot}%
\usepackage{booktabs}%
\usepackage{geometry}%

\raggedbottom

\begin{document}

\title[Land-use footprint of Brazilian soy]{Mapping the land-use footprint
of Brazilian soy embodied in international consumption: a spatially explicit
input--output approach}

\author*[1]{\fnm{Liza} \sur{Buriak}}\email{TODO@wu.ac.at}

\affil*[1]{\orgdiv{Institute for Ecological Economics},
\orgname{Vienna University of Economics and Business (WU)},
\orgaddress{\city{Vienna}, \country{Austria}}}

\abstract{\emph{TODO --- Data-only working document. The full abstract lives in
the main manuscript. Placeholder so the file compiles standalone.}}

\keywords{soy, land footprint, telecoupling, MRIO, FABIO, Brazil, deforestation}

\maketitle

%%%%%%%%%%%%%%%%%%%%%%%%%%%%%%%%%%%%%%%%%%%%%%%%%%%%%%%%%%%%%%%%%%%%%%
\section{Data}\label{sec:data}
%%%%%%%%%%%%%%%%%%%%%%%%%%%%%%%%%%%%%%%%%%%%%%%%%%%%%%%%%%%%%%%%%%%%%%

The model is built entirely from publicly and freely available data. Inputs
fall into four groups: (i) subnational Brazilian statistics on soy production,
processing, domestic use and trade; (ii) spatial layers describing the
multimodal transport network; (iii) administrative and environmental
geographies used to place and attribute the footprint; and (iv) the global
input--output backends (FABIO and EXIOBASE) plus the independent supply-chain
benchmark (Trase). Unless noted otherwise, all series are annual and cover
2010--2020 at the Brazilian municipality (\emph{munic\'ipio}) level
($\sim$5{,}570 units). Where a dataset is provided only for a single reference
year (e.g.\ gridded livestock, household consumption), that year is stated and
the value is held constant, as detailed below. All source URLs were last
accessed in July 2026. A programmatic downloader for the machine-retrievable
sources is provided with the code (\texttt{code/download\_data.py}).

% ==================================================================
\subsection{Subnational production, processing, consumption and trade}
\label{subsec:brazil}
% ==================================================================

\subsubsection{Agricultural production and population (IBGE)}
Municipal soybean production, planted area and harvested area come from the
Brazilian Institute of Geography and Statistics (IBGE) automated recovery
system SIDRA, table 1612 (\emph{Produ\c{c}\~ao Agr\'icola Municipal}),
variables 109 (planted area, ha), 216 (harvested area, ha) and 214 (production,
tonnes), for the soybean class \texttt{c81/2713}. Data are pulled from the
SIDRA API
(\url{https://apisidra.ibge.gov.br/values/t/1612/n6/all/v/109,216,214/p/{year}/c81/2713};
table page \url{https://sidra.ibge.gov.br/tabela/1612}). Resident population by
municipality is taken from SIDRA table 6579 (variable 9324), which runs through
2021; for 2022 the 2022 Census table 9514 is used. Population enters as a
proxy for household soy-oil consumption (Methods).

\subsubsection{Livestock herds (IBGE)}
Municipal herd sizes for ten animal categories (cattle, buffalo, horses, pigs
and breeding sows, goats, sheep, chickens and laying hens, quail) are from
SIDRA table 3939 (\emph{Efetivo dos rebanhos}, variable 105, herd-type
classification \texttt{c79}). Milked-cow counts, used to split dairy from meat
herds, are from SIDRA table 94 (variable 106). Feedlot cattle come from the
IBGE Agricultural Census: table 919 (2006 Census) extrapolated with the ABIEC
growth rate, and table 6911 (2017 Census). All are retrieved from the same
SIDRA API endpoint pattern
(\url{https://apisidra.ibge.gov.br/values/t/{table}/n6/all/...}).

\subsubsection{Soy crushing and refining capacity (ABIOVE)}
Installed processing capacity --- crushing and refining/bottling facilities
with active-facility counts and daily capacities by municipality --- is from
the Brazilian Association of Vegetable Oil Industries (ABIOVE), ``Capacidade
Instalada'' statistics (\url{https://abiove.org.br/estatisticas/}). These
spreadsheets are downloaded manually and stored as
\texttt{Processing\_facilities\_\{year\}\_ABIOVE.xlsx} (sheets
\texttt{processing\_MUN} and \texttt{refining\_bottling\_MUN}).

\subsubsection{Biodiesel capacity (ANP)}
Biodiesel plant capacity (m$^3$/day per municipality) and the regional share of
soy as biodiesel feedstock are from the National Agency of Petroleum, Natural
Gas and Biofuels (ANP) Statistical Yearbook, Table 2.6
(\url{https://www.gov.br/anp/pt-br/centrais-de-conteudo/publicacoes/anuario-estatistico/}),
stored as \texttt{Biodiesel\_capacity\_\{year\}\_ANP.xlsx}.

\subsubsection{Household soy-oil consumption and grain storage (IBGE)}
Per-capita soy-oil acquisition by state, used to distribute national food-oil
use across municipalities, is from the IBGE Household Budget Survey (POF)
2017--2018 --- the most recent available --- held constant across years
(\url{https://www.ibge.gov.br/estatisticas/sociais/educacao/9050-pesquisa-de-orcamentos-familiares.html}).
Grain-storage capacity by municipality, used to allocate stock changes, is from
the IBGE Stocks Survey (\emph{Pesquisa de Estoques}), SIDRA table 278
(\url{https://sidra.ibge.gov.br/tabela/278}).

\subsubsection{National commodity balances (FAO)}
National supply--use balances that the municipal totals are reconciled against
are the FAO Food/Commodity Balance Sheets for Brazil (area code 21), items 2555
(soybeans), 2571 (soybean oil) and 2590 (soybean cake), covering production,
imports, exports, food, feed, seed, processing, other uses and stock variation.
These are extracted from the FAOSTAT bulk download
(\url{https://bulks-faostat.fao.org/production/FoodBalanceSheets_E_All_Data_(Normalized).zip};
domain page \url{https://www.fao.org/faostat/en/#data/FBS}).

\subsubsection{Livestock distribution and feed rates (FAO)}
Herds are disaggregated into production systems using the FAO Gridded Livestock
of the World (GLW3, 2010 reference year) density rasters for cattle, buffalo,
pigs (extensive/intensive/industrial) and chickens (extensive/intensive)
\cite{gilbert2018glw}, distributed via the Harvard Dataverse
(\url{https://dataverse.harvard.edu/dataverse/glw}), together with the FAO
GLEAM ruminant production-system classification
(\texttt{glps\_gleam\_61113\_10km.tif}) and GLEAM per-animal feed ratios
(dry-matter intake and soy/soy-cake shares) \cite{fao2018gleam},
\url{https://www.fao.org/gleam/resources/en/}. These are static 2010 layers,
reused for all years; the resulting production-system \emph{shares} (not herd
sizes, which are annual) are therefore held constant over the study period.

\subsubsection{Municipal-level trade (COMEX/MDIC)}
Municipality-level exports and imports are from the Brazilian foreign-trade
database COMEX Stat (Ministry of Development, Industry, Trade and Services),
filtered to the soy HS4 codes 1201 (soybeans), 1507 (soybean oil) and 2304
(soybean cake). Per-year municipal files and the country and municipality
lookup tables are downloaded from
\url{https://balanca.economia.gov.br/balanca/bd/comexstat-bd/mun/}
(files \texttt{EXP\_\{year\}\_MUN.csv}, \texttt{IMP\_\{year\}\_MUN.csv}) and
\url{https://balanca.economia.gov.br/balanca/bd/tabelas/}
(\texttt{PAIS.csv}, \texttt{UF\_MUN.csv}).

% ==================================================================
\subsection{Transport network}\label{subsec:transport-data}
% ==================================================================
The subnational routing of physical flows (Methods) uses four multimodal
network layers:
\begin{itemize}
  \item \textbf{Roads.} OpenStreetMap road network
    (\texttt{gis\_osm\_roads\_free\_1.shp}), obtained as a Brazil extract from
    Geofabrik (\url{https://download.geofabrik.de/south-america/brazil.html});
    \textcopyright{} OpenStreetMap contributors.
  \item \textbf{Waterways.} Federal inland-waterway corridors
    (\emph{Hidrovias}) from the National Department of Transport Infrastructure
    (DNIT), \url{https://www.gov.br/dnit/pt-br}.
  \item \textbf{Railways.} Rail lines, stations and annual rail-cargo flows by
    commodity (soybean, soy cake, vegetable oil) from the National Land
    Transport Agency (ANTT), 2006--2021
    (\url{https://dados.antt.gov.br/dataset/transporte-ferroviario-de-cargas-e-passageiros}).
  \item \textbf{Ports and port cargo.} Port locations and annual waterway-cargo
    records from the National Waterway Transport Agency (ANTAQ) statistical
    system, 2017--2022 with a 2013 fallback for earlier years
    (\url{https://web.antaq.gov.br/anuario/}).
\end{itemize}
A small number of soy-relevant rail stations and supplementary port points were
hand-curated and are distributed with the code.

% ==================================================================
\subsection{Administrative and environmental geographies}\label{subsec:geo}
% ==================================================================
\textbf{Boundaries.} Municipal polygons are IBGE administrative meshes,
retrieved with the \texttt{geobr} R package \cite{pereira2019geobr}
(\url{https://ipeagit.github.io/geobr/}; also on CRAN,
\url{https://cran.r-project.org/package=geobr}) and municipality code/name
lookups from the IBGE localities API
(\url{https://servicodados.ibge.gov.br/api/v1/localidades/municipios});
state-level boundaries for mapping use GADM v3.6
(\texttt{gadm36\_BRA\_1}, \url{https://gadm.org/}).

\textbf{Biomes and frontier.} Footprint origin is attributed to biomes using
the IBGE 1:250{,}000 biome layer (\texttt{lm\_bioma\_250.shp}: Amaz\^onia,
Cerrado, Caatinga, Mata Atl\^antica, Pantanal, Pampa;
\url{https://www.ibge.gov.br/geociencias/informacoes-ambientais/estudos-ambientais/15842-biomas.html}).
The agricultural \emph{frontier} is defined as the Amazon biome plus MATOPIBA,
the Cerrado expansion zone spanning Maranh\~ao, Tocantins, Piau\'i and Bahia.
MATOPIBA is operationalized as the intersection of the Cerrado biome with these
four states, which closely approximates the official Embrapa delimitation of
the region \cite{miranda2014matopiba}.

\textbf{Grid refinement (optional).} An optional 30\,m refinement of the
municipal footprint uses MapBiomas annual soy land-use tiles
\cite{souza2020mapbiomas} (\url{https://mapbiomas.org/}), exported from Google
Earth Engine with an
adapted version of the official MapBiomas user toolkit
(\url{https://github.com/mapbiomas-brazil/user-toolkit}). This layer is a
visual refinement and is not load-bearing for the headline results.

% ==================================================================
\subsection{Global supply-chain models and benchmark}\label{subsec:mrio-data}
% ==================================================================
\textbf{FABIO.} Subnational flows are nested into the Food and Agriculture
Biomass Input--Output model, version~2 (FABIO v2), built by the fineprint team
(\url{https://github.com/fineprint-global/fabio}). The v2 matrices cover
2010--2023 at a resolution of 22{,}263 processes with 123 environmental
extensions (files \texttt{E}, \texttt{Z\_mass\_b}, \texttt{Z\_value\_b},
\texttt{Y\_b}), complemented by the v2 build ingredients (balanced commodity
balance sheets, bilateral trade, tidied livestock and price tables) at the v2
classification (181 regions $\times$ 123 items). The published FABIO v1.1
release (1986--2013; DOI \texttt{10.5281/zenodo.2577067}) is cited for
provenance only and is \emph{not} used, as its classification and time span
differ \cite{bruckner2019fabio}.

\textbf{EXIOBASE.} The non-soy ``background'' of the economy is supplied by
EXIOBASE~3 \cite{stadler2018exiobase} in the product-by-product (pxp)
transformation, 2000--2020, with 9{,}800 $\times$ 9{,}800 flow matrices (49
regions $\times$ 200 products) and their environmental satellites
(\url{https://www.exiobase.eu/}; data on Zenodo, DOI
\texttt{10.5281/zenodo.3583070}). These are hybridised with FABIO for the
non-soy quadrant of the footprint model.

\textbf{Bilateral trade (FAO).} Re-export disentangling and the FABIO trade
layer draw on the FAOSTAT Detailed Trade Matrix for Brazil (items 236
soybeans, 237 soybean oil, 238 soybean cake), from the FAOSTAT bulk download
(\url{https://bulks-faostat.fao.org/production/Trade_DetailedTradeMatrix_E_All_Data_(Normalized).zip}).

\textbf{Independent benchmark (Trase).} Model flows from origin municipality to
country of first import are validated against the Trase Brazil soy composite
dataset, version 2.6.1 (\texttt{brazil\_soy\_v2\_6\_1\_composite.csv},
2004--2022), from the Trase initiative (\url{https://trase.earth/})
\cite{godar2015sei,escobar2020trase}. Trase is used exclusively as a
validation benchmark; no Trase data enter the model itself.

% ==================================================================
\subsection{Coverage and summary}\label{subsec:coverage}
% ==================================================================
The binding constraint on temporal coverage is FABIO v2 (2010--2023), so the
footprint series is reported for the decade 2010--2020; the origin-to-importer
flow benchmark against Trase extends back to 2004 and supports a longer
validation backdrop. Table~\ref{tab:data} summarises all inputs.

\begin{table}[htbp]
\small
\caption{Input datasets, by role. All are publicly and freely available;
per-source access details and URLs are given in the text. ``Munis''
$=$ Brazilian municipalities.}\label{tab:data}
\begin{tabular}{@{}p{3.5cm}p{2.3cm}p{3.0cm}p{2.6cm}@{}}
\toprule
Dataset & Provider & Role & Coverage \\
\midrule
Soy production, area (t.1612)      & IBGE (SIDRA)        & production            & munis, annual \\
Population (t.6579/9514)           & IBGE (SIDRA)        & food-use proxy        & munis, annual \\
Herds (t.3939), milk cows (t.94)   & IBGE (SIDRA)        & feed demand           & munis, annual \\
Feedlot cattle (t.919/6911)        & IBGE (Census)       & feed demand           & munis, 2006/2017 \\
Crush/refining capacity            & ABIOVE              & processing            & munis, annual \\
Biodiesel capacity                 & ANP                 & other (biodiesel) use & munis, annual \\
Soy-oil acquisition (POF)          & IBGE                & food-oil allocation   & states, 2017--18 \\
Grain storage (t.278)              & IBGE                & stock allocation      & munis, annual \\
Commodity balance (FBS)            & FAO (FAOSTAT)       & national reconciliation & Brazil, annual \\
Gridded livestock (GLW3)           & FAO (Dataverse)     & system split          & 10\,km, 2010 \\
Production systems, feed (GLEAM)    & FAO                 & feed rates            & 10\,km, static \\
Municipal trade                    & COMEX/MDIC          & exports/imports       & munis, annual \\
Road network                       & OSM (Geofabrik)     & transport             & national \\
Waterways                          & DNIT                & transport             & national \\
Rail lines, stations, cargo        & ANTT                & transport             & national, 2006--21 \\
Ports, port cargo                  & ANTAQ               & transport             & national, 2017--22 \\
Municipal boundaries               & IBGE (geobr)        & geometry              & munis \\
State boundaries                   & GADM v3.6           & mapping               & states \\
Biomes 1:250k                      & IBGE                & origin attribution    & national \\
Soy land tiles                     & MapBiomas (GEE)     & grid refinement       & 30\,m, annual \\
FABIO v2 (MRIO)                    & fineprint           & supply-chain nesting  & global, 2010--23 \\
EXIOBASE 3 pxp                     & EXIOBASE            & non-soy background    & global, 2000--20 \\
Detailed trade matrix              & FAO (FAOSTAT)       & re-exports/trade      & Brazil, annual \\
Trase composite v2.6.1             & Trase               & validation benchmark  & munis, 2004--22 \\
\botrule
\end{tabular}
\end{table}

% ==================================================================
\subsection{Data availability}\label{subsec:data-availability}
% ==================================================================
All input data used in this study are publicly and freely available from the
providers listed above. They are not redistributed with the code because of
differing licences; a download script with direct source links is provided
(\texttt{code/download\_data.py}), and the datasets requiring manual retrieval
(ABIOVE, ANP, POF, MapBiomas tiles) are documented with step-by-step
instructions in the repository. The intermediate and final model outputs
underlying the figures will be deposited on Zenodo upon publication.

% ==================================================================
\subsection{Code availability}\label{subsec:code-availability}
% ==================================================================
The complete modelling pipeline (data preparation, transport allocation,
FABIO/EXIOBASE nesting, footprint accounting, validation and figures) is
written in R and Python and released under the GNU GPL~v3 at
\texttt{<repository URL --- TODO>}. The results reported here have no
proprietary software dependency. An experimental multimodal-transport module
included in the repository (not used for the reported results) additionally
requires a GAMS licence.

%%%%%%%%%%%%%%%%%%%%%%%%%%%%%%%%%%%%%%%%%%%%%%%%%%%%%%%%%%%%%%%%%%%%%%
\bibliography{refs}%
%%%%%%%%%%%%%%%%%%%%%%%%%%%%%%%%%%%%%%%%%%%%%%%%%%%%%%%%%%%%%%%%%%%%%%

\end{document}
